\newcommand{\wineIntro}{But we also reflect on what ``freedom'' means to different people and what it means to each of us.}

\newcommand{\wineSummaryThisYear}{This year, we think about different views of freedom.}

\newcommand{\wineFirstCup}{
    \newReading With this first cup of wine, we compare how some individuals define their own view of freedom.

    

    \subsubsection*{What is freedom?}
    What is freedom? How is freedom related to liberty? Liberty conveys the idea of the right or the ability of the individual to do what he wants. Freedom on the other hand is a word which contains inbuilt restrictions. 

    My freedom to swing my hands stops at my neighbor's nose. A free society contains individuals who are free to do what they do, whilst respecting the freedoms of others. Freedom is the liberty of the individual coupled with concern by the liberated individual for the liberties of others.

    Freedom cannot be defined, except through an analysis of the restrictions on human action.
    \textemdash{}Dr. Mark Cooray

    \subsubsection*{``Freedom's just another word for nothing' left to lose''}
    \newReading Sure, having nothing to lose is a type of freedom. But it's a temporary type; freedom is certainly not synonymous with having nothing to lose. And if you really, truly have nothing at all to lose, you're not even remotely free; you're closer to slavery than freedom, since you'd better get cracking to find yourself something to eat within the next 24 hours or so.

    Real freedom, it seems to me, comes when you don't have to do what others tell you. Granted, possessions always bring responsibilities, but carefully chosen possessions can provide more freedom than they take away.

    I've always thought a better statement is: ``Freedom means never having to say `I'm sorry' to someone who's really a total jerk.''
    
    \textemdash{}Posted by ``jackelope'' on the StraightDope message boards 

    \subsubsection*{Free to be yourself}
    \newReading The most important kind of freedom is to be what you really are. You trade in your reality for a role. You trade in your senses for an act. You give up your ability to feel and in exchange, put on a mask. I think people resist freedom because they're afraid of the unknown.

    The only solution is to confront yourself with the greatest fear imaginable. Expose yourself to yourself and to your deepest fear. After that, fear has no power, and fear of freedom shrinks and vanishes. You are free.

    \textemdash{}Jim Morrison

    With this first cup of wine, we begin to compose our own definition of freedom. 


}

\newcommand{\wineSecondCup}{
    \newReading With this second cup of wine, we contemplate different ways to view freedom in our lifestyles.

    \subsubsection*{The Monk and the Minister}
    Two close boyhood friends grow up and go their separate ways.  One becomes a humble monk, the other a rich and powerful minister to the king.

    Years later they meet.  As they catch up, the minister (in his fine robes) takes pity on the thin, shabby monk.  Seeking to help, he says: ``You know, if you could learn to cater to the king you wouldn't have to live on rice and beans.''

    To which the monk replies: ``If you could learn to live on rice and beans, you wouldn't have to cater to the king.''

    \newReading While most all of us fall somewhere between these two examples, they provide a spectrum on which to consider freedom. 

    Forming business deals or other contracts, and making family or personal promises are actions we take freely that can bring great happiness, but at the same time they limit our freedom to follow other pursuits by binding us to those obligations. 

    For this second cup of wine, we ask ourselves: Is my view of freedom closer to the monk, or to the minister? And are my choices and actions leading me towards or away from that goal?
}

\newcommand{\wineThirdCup}{
    \newReading The third cup of wine symbolizes spiritual freedom. Our people have known the need for spiritual resistance in many ages. Even in the worst of circumstances, we have maintained our dignity. Even in the concentration camps, many Jews scrupulously observed \textit{Halakhah} (Jewish law) in defiance of the oppression they suffered at the hands of the Nazis. 
    
    \newReading One such Jew addressed this question to Rabbi Ephraim Oshry: ``Should a Jew, having to perform forced labor for the Nazis, continue to recite the benediction in the morning prayers: We praise You, Adonai our God, ruler of the universe, who has not made me a slave?'' Rabbi Oshry responded: ``Heaven forbid that we should give up reciting this blessing. On the contrary, now of all times we are obliged to say this blessing so that our adversaries and tormentors realize that, although we are in their power to do with us as their wicked machinations devise, we nonetheless perceive ourselves not as slaves, but as free people, prisoners for the time being, whose liberation will soon come.''
    
    With this third cup of wine, let us seek the spiritual freedom that generations before us sacrificed to maintain. Let us open our hearts and minds to experience God in our own lives.
}

\newcommand{\wineFourthCup}{
    \newReading As our seder draws near to the end, we take up our cups one last time. The redemption is not yet complete. Not everyone in our world is yet free. This fourth cup reminds us of our responsibility to be God's partners in bringing freedom to those enslaved, peace to those at war, food to those who hunger. This is our purpose as Jews. May we live to fulfill it.
}